\chapter{Low Vision}
\index{Low Vision}
\label{sec:Low Vision}

\section{Overview}
\begin{itemize}[noitemsep]
  \item Low vision is a significant, uncorrectable loss of vision that interferes with daily activities, defined as best-corrected visual acuity of 20/70 or worse, or a visual field of 20 degrees or less
  \item It is not blindness; people with low vision retain useful vision and can be helped with rehabilitation
  \item It affects millions of older adults worldwide and is most often caused by age-related eye disease
\end{itemize}
\section{Causes}
Low vision happens because of conditions that cannot be fully corrected with glasses, surgery, or medication. Common causes include age-related macular degeneration, glaucoma, diabetic retinopathy, cataracts, and retinal disease such as retinitis pigmentosa or retinal vein occlusion. Stroke affecting the visual pathways can also cause low vision.
\section{Risk Factors}
\begin{itemize}[noitemsep]
  \item Age over 65 years
  \item Diabetes and high blood pressure
  \item Smoking
  \item Family history of eye disease
  \item Uncorrected refractive error over long periods
  \item Eye trauma or prior eye surgery complications
\end{itemize}
\section{Signs and Symptoms}
\subsection{Common symptoms}
\begin{itemize}[noitemsep]
  \item Blurred or distorted vision that glasses cannot fix
  \item Difficulty reading, writing, driving, or recognizing faces
  \item Loss of peripheral (side) vision
  \item Difficulty with glare and low-light conditions
  \item Central blind spots or tunnel vision
  \item Impaired depth perception and increased risk of falls
\end{itemize}
\subsection{Emergency warning signs}
\begin{itemize}[noitemsep]
  \item Sudden or rapidly worsening vision loss requires immediate evaluation
\end{itemize}
\section{Progression and Stages}
Low vision may be stable, progressive, or episodic depending on the underlying cause. Age-related macular degeneration and glaucoma tend to progress gradually, while stroke or vein occlusion may cause sudden loss. The functional impact depends on the residual vision, the specific tasks affected, and the effectiveness of rehabilitation.
\section{Complications}
\begin{itemize}[noitemsep]
  \item Increased risk of falls, fractures, and injury
  \item Difficulty managing medications and daily tasks
  \item Depression, social isolation, and reduced independence
  \item Loss of driving ability
  \item Decreased quality of life
\end{itemize}
\section{Diagnosis}
\begin{itemize}[noitemsep]
  \item Comprehensive eye examination including visual acuity and visual field testing
  \item Low-vision assessment to determine functional needs
  \item Optical coherence tomography (OCT), fundus photography, and other imaging depending on the cause
  \item Evaluation of the underlying disease by an ophthalmologist
  \item Referral to low-vision and rehabilitation services
\end{itemize}
\section{Prevention}
\begin{itemize}[noitemsep]
  \item Early treatment of the underlying eye disease
  \item Control of diabetes, blood pressure, and cholesterol
  \item Regular comprehensive eye examinations, especially after age 50
  \item Quitting smoking and maintaining a healthy lifestyle
  \item Fall-prevention measures in the home
\end{itemize}
\section{Treatment Options}
\begin{itemize}[noitemsep]
  \item Low-vision aids: magnifiers, telescopic lenses, high-contrast and large-print materials
  \item Electronic aids: screen readers, video magnifiers (CCTV), speech-to-text software
  \item Occupational therapy and low-vision rehabilitation
  \item Lighting and contrast optimization in the home
  \item Mobility training and orientation training for severe cases
  \item Treatment of the underlying eye disease where possible
\end{itemize}
\section{Living with It}
\begin{itemize}[noitemsep]
  \item Work with a low-vision specialist to find the right aids
  \item Improve home lighting, contrast, and organization
  \item Seek emotional support and connect with low-vision support groups
  \item Continue participating in hobbies with adapted techniques
  \item Consider counseling for depression or anxiety
\end{itemize}
\section{Outlook}
The outlook depends on the underlying cause, but most people with low vision can maintain independence and good quality of life with rehabilitation and assistive technology. Progression varies with the specific disease; early and ongoing support significantly improves outcomes.
